\chaptertitle{Bibliography}

                                     Bibliography
  The presence of two asterisks indicates that the hook or article was a prime
  motivator of my book. The presence of a single asterisk means that
  the book or article has some special feature or quirk which I want to single out.
          I have not given many direct pointers into technical literature; instead I have
  chosen to give "meta-pointers": pointers to books which have pointers to technical
  literature.

Allen, John. The Anatomy of LISP. New York: McGraw-Hill, 1978. The most
 comprehensive book on LISP, the computer language which has dominated Artificial
 Intelligence research for two decades. Clear and crisp.
Anderson, Alan Ross, ed. Minds and Machines. Englewood Cliffs, N. J.: PrenticeHall,
 1964. Paperback. A collection of provocative articles for and against Artificial
 Intelligence. Included are Turing's famous article "Computing Machines and
 Intelligence" and Lucas' exasperating article "Minds, Machines, and Godel".
Babbage, Charles. Passages from the Life of a Philosopher. London: Longman, Green,
 1864. Reprinted in 1968 by Dawsons of Pall Mall (London). A rambling selection of
 events and musings in the life of this little-understood genius. There's even a play
 starring Turnstile, a retired philosopher turned politician, whose favorite musical
 instrument is the barrel-organ. I find it quite jolly reading.
Baker, Adolph. Modern Physics and Anti-physics. Reading, Mass.: Addison-Wesley,
 1970. Paperback. A book on modern physics-especially quantum mechanics and
 relativity-whose unusual feature is a set of dialogues between a "Poet" (an antiscience
 "freak") and a "Physicist". These dialogues illustrate the strange problems which arise
 when one person uses logical thinking in defense of itself while another turns logic
 against itself.
Ball, W. W. Rouse. "Calculating Prodigies", in James R. Newman, ed. The World of
 Mathematics, Vol. 1. New York: Simon and Schuster, 1956. Intriguing descriptions of
 several different people with amazing abilities that rival computing machines.
Barker, Stephen F. Philosophy of Mathematics. Englewood Cliffs, N. J.: Prentice-Hall.
 1969. A short paperback which discusses Euclidean and non-Euclidean geometry, and
 then Godel's Theorem and related results without any mathematical forms lism.
* Beckmann, Petr. A History of Pi. New York: St. Martin's Press, 1976. Paperback.
 Actually, a history of the world, with pi as its focus. Most entertaining, as well as a
 useful reference on the history of mathematics.
* Bell, Eric Temple. Men of Mathematics. New York: Simon \& Schuster, 1965.
 Paperback. Perhaps the most romantic writer of all time on the history of mathematics.
 He makes every life story read like a short novel. Nonmathematicians can come away
 with a true sense of the power, beauty, and meaning of mathematics.
Benacerraf, Paul. "God, the Devil, and Godel". Monist 51 (1967): 9. One of the most
 important of the many attempts at refutation of Lucas. All about mechanism and
 metaphysics, in the light of Godel's work.
Benacerraf, Paul, and Hilary Putnam. Philosophy of Mathematics-Selected Readings.
 Englewood Cliffs, N. J.: Prentice-Hall, 1964. Articles by Godel, Russell, Nagel, von

Bibliography                                                                            745

  Neumann, Brouwer, Frege, Hilbert, Poincare, Wittgenstein, Carnap, Quine, and others
  on the reality of numbers and sets, the nature of mathematical truth, and so on.
* Bergerson, Howard. Palindromes and Anagrams. New York: Dover Publications, 1973.
  Paperback. An incredible collection of some of the most bizarre and unbelievable
  wordplay in English. Palindromic poems, plays, stories, and so on.
Bobrow, D. G., and Allan Collins, eds. Representation and Understanding: Studies in
  Cognitive Science. New York: Academic Press, 1975. Various experts on Artificial
  Intelligence thrash about, debating the nature of the elusive "frames", the question of
  procedural vs. declarative representation of knowledge, and so on. In a way, this book
  marks the start of a new era of Al: the era of representation.
* Boden, Margaret. Artificial Intelligence and Natural Man. New York: Basic Books,
  1977. The best hook I have ever seen on nearly all aspects of Artificial Intelligence,
  including technical questions, philosophical questions, etc. It is a rich book, and in niv
  opinion, a classic. Continues the British tradition of clear thinking and expression on
  matters of mind, free will, etc. Also contains an extensive technical bibliography.
. Purposive Explanation in Psychology. Cambridge, Mass.: Harvard University Press,
  1972. The book to which her Al book is merely- "an extended footnote", says Boden.
* Boeke, Kees. Cosmic View: The Universe in 40 Jumps. New York: John Day, 1957.
  The ultimate book on levels of description. Everyone should see this book at some
  point in their life. Suitable for children.
** Bongard, M. Pattern Recognition. Rochelle Park, N. J.: Hay den Book Co., Spartan
  Books, 1970. The author is concerned with problems of determining categories in an
  ill-defined space. In his book, he sets forth a magnificent collection of 100 "Bongard
  problems" (as I call them)-puzzles for a pattern recognizer (human or machine) to test
  its wits on. They are invaluably stimulating for anyone whoo is interested in the nature
  of intelligence.
Boolos, George S., and Richard Jeffrey. Computability and Logic. New York: Cambridge
  University Press. 1974. A sequel to Jeffrey's Formal Logic. It contains a wide number
  of results not easily obtainable elsewhere. Quite rigorous, but this does not impair its
  readability.
Carroll, John B., Peter Davies, and Barry Rickman. The American Heritage Word
  Frequency Book. Boston: Houghton Mifflin, and New York: American Heritage
  Publishing Co., 1971. A table of words in order of frequency in modern written
  American English. Perusing it reseals fascinating things about out- thought processes.
Cerf, Vinton. "Parry Encounters the Doctor". Datamation, July 1973, pp. 62-64. The first
  meeting of artificial "minds"-what a shock!
Chadwick, John. The Decipherment of Linear B. New York: Cambridge University
  Press, 1958. Paperback. A book about a classic decipherment-that of a script from the
  island of Crete-done by a single man: Michael Ventris.
Chaitin, Gregory J. "Randomness and Mathematical Proof"'. Scientific American, May
  1975. An article about an algorithmic definition of randomness, and its intimate relation
  to simplicity. These two concepts are tied in with Godel's Theorem, which assumes a
  new meaning. An important article.
Cohen, Paul C. Set Theory and the Continuum Hypothesis. Menlo Park, Calif.: W. A.
  Benjamin, 1966. Paperback. A great contribution to modern mathematics-the
  demonstration that various statements are undecidable within the usual formalisms for

Bibliography                                                                            746

  set theory-is here explained to nonspecialists by its discoverer. The necessary
  prerequisites in mathematical logic are quickly, concisely, and quite clearly presented.
Cooke, Deryck. The Language of Music. New York: Oxford University Press, 1959.
  Paperback. The only book that I know which tries to draw an explicit connection
  between elements of music and elements of human emotion. A valuable start down
  what is sure to be a long hard road to understanding music and the human mind.
* David, Hans Theodore.J. S. Bach's Musical Offering. New York: Dover Publications,
  1972. Paperback. Subtitled "History, Interpretation, and Analysis". A wealth of
  information about this tour de force by Bach. Attractively written.
** David, Hans Theodore, and Arthur Mendel. The Bach Reader. New York: W. W.
  Norton, 1966. Paperback. An excellent annotated collection of original source material
  on Bach's life, containing pictures, reproductions of manuscript pages, many short
  quotes from contemporaries, anecdotes, etc., etc.
Davis, Martin. The Undecidable. Hewlett, N. Y.: Raven Press, 1965. An anthology of
  some of the most important papers in metamathematics from 1931 onwards (thus quite
  complementary to van Heijenoort's anthology). Included are a translation of Godel's
  1931 paper, lecture notes from a course which Godel once gave on his results, and then
  papers by Church, Kleene, Rosser, Post, and Turing.
Davis, Martin, and Reuben Hersh. "Hilbert's Tenth Problem". Scientific American,
  November 1973, p. 84. How a famous problem in number theory was finally shown to
  be unsolvable, by a twenty-two-year old Russian.
** DeLong, Howard. A Profile of Mathematical Logic. Reading, Mass.: Addison-
  Wesley, 1970. An extremely carefully written book about mathematical logic, with an
  exposition of Godel's Theorem and discussions of many philosophical questions. One
  of its strong features is its outstanding, fully annotated bibliography. A book which
  influenced me greatly.
Doblhofer, Ernst. Voices in Stone. New York: Macmillan, Collier Books, 1961. Paper
  back. A good hook on the decipherment of ancient scripts.
* Drevfus, Hubert. What Computers Can't Do: A Critique of Artificial Reason. New
  York:Harper \& Roc-, 1972. A collection of many arguments against Artificial
  Intelligence from someone outside of the field. Interesting to try to refute. The Al
  community and Drevfus enjoy a relation of strong mutual antagonism. It is important to
  have people like Dreyfus around, even if you find them very irritating.
Edwards, Harold M. "Fermat's Last Theorem". Scientific American, October 1978, pp.
  104-122. A complete discussion of this hardest of all mathematical nuts to crack, from
         its origins to the most modern results. Excellently illustrated.
* Ernst, Bruno. The Magic Mirror of M. C. Escher. New York: Random House, 1976.
  Paperback. Escher as a human being, and the origins of his drawings, are discussed
  with devotion by a friend of mans' years. A "must" for any lover of Escher.
** Escher, Maurits C., et al. The World of M. C. Escher. New York: Harry N. Abrarns,
1972. Paperback. The most extensive collection of reproductions of Escher's works.
  Escher comes about as close as one can to recursion in art, and captures the spirit of
  Godel's Theorem in some of his drawings amazingly well,
Feigenbaum, Edward, and Julian Feldman, eds. Computers and Thought. New York:
  McGraw-Hill, 1963. Although it is a little old now, this book is still an important
  collection of ideas about Artificial Intelligence. Included are articles on Gelernters

Bibliography                                                                          747

  geometry program, Samuel's checkers program, and others on pattern recognition,
  language understanding, philosophy, and so on.
Finsler, Paul. "Formal Proofs and Undecidability", Reprinted in van Heijenoorcs
  anthology From Frege to Godel (see below). A forerunner of Godel's paper, in which
  the existence of undecidable mathematical statements is suggested, though not
  rigorously demonstrated.
Fitzpatrick, P. J. "To Godel via Babel"- Mind 75 (1966): 332-350. An innovative
  exposition of Godel's proof which distinguishes between the relevant levels by using
  three different languages: English, French, and Latin'
* Gablik, Suzi. Magritte. Boston, Mass.: New York Graphic Society, 1976. Paperback.
  An excellent book on Magritte and his works by someone who really understands their
  setting in a wide sense; has a good selection of reproductions.
* Gardner, Martin. Fads and Fallacies. New York: Dover Publications, 1952. Paperback.
  Still probably the best of all the anti-occult books. Although probably not intended as a
  book on the philosophy of science, this book contains many lessons therein. Over and
  over, one faces the question, "What is evidence?" Gardner demonstrates how
  unearthing "the truth" requires art as much as science.
Gebstadter, Egbert B. Copper, Silver, Gold: an Indestructible Metallic Alloy. Perth:
  Acidic Books, 1979. A formidable hodge-podge, turgid and confused-yet remarkably
  similar to the present work. Professor Gebstadter's Shandean digressions include some
  excellent examples of indirect self-reference. Of particular interest is a reference in its
  well-annotated bibliography to an isomorphic, but imaginary, book.
** Godel, Kurt. On Formally Undecidable Propositions. New York: Basic Books, 1962.
  A translation of Godel's 1931 paper, together with some discussion.
. "Uber Formal Unentscheidbare Satze der Principia Mathematica and Verwandter
  Systeme, I." Monatshefte fur Mathematik and Physik, 38 (1931), 173-198. Godel's
  1931 paper.
Golf man, Ering. Frame Analysis. New York: Harper \& Row, Colophon Books, 1974.
  Paperback. A long documentation of the definition of "systems" in human
  communication, and how in art and advertising and feporting and the theatre, the
  borderline between "the system" and "the world" is perceived and exploited and
  violated.
Goldstein, Ira. and Seymour Papert. "Artificial Intelligence, Language, and the Study of
  Knowledge". Cognitive Science 1 (January l9/ 7): 84-123. A survey article concerned
  with the past and future of Al. The authors see three periods so far: "Classic",
  "Romantic,., and "Modern".
Guod, I. I. "Human and Machine Logic". British Journal for the Philosophy of Science In
  (196/): 144. One of the most interesting attempts to refute Lucas, having to do with
  whether the repeated application of the diagonal method is itself a mechanizable
  operation.
. "Godel's Theorem is a Red Herring". British Journal for the Philosophy of Science 19
  (1969): 357. In which Good maintains that Lucas' argument has nothing to do with
  Godel's Theorem, and that Lucas should in fact have entitled his article "Minds,
  Machines, and Transfinite Counting". The Good-Lucas repartee is fascinating.
Goodman, Nelson. Fact, Fiction, and Forecast. 3rd ed. Indianapolis: Bobbs-Merrill, 1973.
  Paperback. A discussion of contrary-to-fact conditionals and inductive logic, including

Bibliography                                                                             748

  Goodman's famous problem-words "bleen" and "grue". Bears very much on the
  question of how humans perceive the world, and therefore interesting especially from
  the Al perspective.
Press, 1976. A short hook about bees, apes, and other animals, and whether or not they
  are "conscious"-and particularh whether or not it is legitimate to use the word
  "consciousness" in scientific explanations of animal behavior.
deGroot, Adriaan. Thought and Choice in Chess. The Hague: Mouton, 1965. A thorough
  study in cognitive psychology, reporting on experiments that have a classical simplicity
  and elegance.
Gunderson, Keith. Mentality and Machines. New York: Doubleday, Anchor Books, 1971.
  Paperback. A very anti-Al person tells why. Sometimes hilarious.
Hanawalt, Philip C., and Robert H. Haynes, eds. The Chemical Basis of Life. San
  Francisco: W. H. Freeman, 1973. Paperback. An excellent collection of reprints from
  the Scientilic American. One of the best ways to get a feeling for what molecular
  biology is about.
Hardy G. H. and E. M. Wright. An Introduction to the Theory of Numbers, 4th ed. New
  York: Oxford University Press, 1960. The classic book on number theory. Chock-full
  of information about those mysterious entities. the whole numbers.
Harmon, Leon. "The Recognition of Faces". Scientific American, November 1973, p. 70.
  Explorations concerning how we represent faces in our memories, and how much
  information is needed in what form for us to be able to recognize a face. One of the
  most fascinating of pattern recognition problems.
van Heijenoort, Jean. From Frege to Godel: A Source Book in Mathematical Logic.
  Cambridge, Mass.: Harvard University Press, 1977. Paperback. A collection of epoch-
  making articles on mathematical logic, all leading up to Godel's climactic revelation,
  which is the final paper in the book.
Henri, Adrian. Total Art: Environments, Happenings, and Performances. New York:
  Praeger, 1974. Paperback. In which it is shown how meaning has degenerated so far m
  modern art that the absence of meaning becomes profoundly meaningful (whatever that
  means).
  * Goodstein, R. L. Development of Mathematical Logic. New York: Springer Verlag,
  1971. A concise survey of mathematical logic, including much material not easily
  found elsewhere. An enjoyable book, and useful as a reference.
Gordon, Cyrus. Forgotten Scripts. New York: Basic Books, 1968. A short and nicely
  written account of the decipherment of ancient hieroglyphics, cuneiform, and other
  scripts.
Griffin, Donald. The Question of Animal Awareness. New York: Rockefeller University
von Foerster, Heinz and James W. Beauchamp, eds. Music by Computers. New York:
        **
John Wiles', 1969. This book contains not only a set of articles about various types of
  computer-produced music, but also a set of four small phonograph records so you can
  actually hear (and judge) the pieces described. Among the pieces is Max Mathews'
  mixture of "Johnny Comes Marching Home" and "The British Grenadiers".
Fraenkel, Abraham, Yehoshua Bar-Hillel, and Azriel Levy. Foundations of Set Theory,
  2nd ed. Atlantic Highlands, N. J.: Humanities Press, 1973. A fairly nontechnical



Bibliography                                                                          749

  discussion of set theory, logic, limitative Theorems and undecidable statements.
  Included is a long treatment of intuitionism.        *
Frev, Peter W. Chess Skill in Man and Machine. New York: Springer Verlag, 1977.
  Anexcellent survey of contemporary ideas in computer chess: why programs work, why
  they don't work, retrospects and prospects.
Friedman, Daniel P. The Little Lisper. Palo Alto, Calif.: Science Research Associates,
  1974. Paperback. An easily digested introduction to recursive thinking in LISP. You'll
  eat it up!
•         Hoare, C. A. R, and D. C. S. Allison. "Incomputability". Computing Surveys 4,
     no. 3 September 1972). A smoothly presented exposition of why the halting problem
     is unsolvable. Proves this fundamental theorem: ''Any language containing
     conditionals and recursive function definitions which is powerful enough to program
     its own interpreter cannot be used to program its own 'terminates' function."
•         Hofstadter, Douglas R. "Energy levels and wave functions of Bloch electrons in
     rational and irrational magnetic fields". Physical Review B, 14, no. 6 (15 September
     1976). The author's Ph.D. work, presented as a paper. Details the origin of "Gplot",
     the recursive graph shown in Figure 34.
Hook, Sidney. ed. Dimensions of Mind. New York: Macmillan, Collier Books, 1961.
  Paperback. A collection of articles on the mind-body problem and the mind-computer
  problem. Some rather strong-minded entries here.
• Hornev, Karen. Self-Analysis. New York: W. W. Norton, 1942. Paperback. A
  fascinating description of how the levels of the self must tangle to grapple with
  problems of self-definition of any individual in this complex world. Humane and
  insightful.
Hubbard, John I. The Biological Basis of Mental Activity. Reading, Mass.:
  AddisonWesley, 1975. Paperback. Just one more book about the brain, with one special
  virtue. however: it contains man long lists of questions for the reader to ponder, and
  references to articles which treat those questions.
• Jackson. Philip C. Introduction to Artificial Intelligence. New York, Petrocelli Charter.
  1975. A recent book, describing, with some exuberance, the ideas of Al. There are a
  huge number of vaguely suggested ideas floating around this book, and for that reason
  it is very stimulating just to page through it. Has a giant bibliography, which is another
  reason to recommend it.
Jacobs, Robert L. Understanding Harmony. New York: Oxford University Press, 1958.
  Paperback. A straightforward book on harmony, which can lead one to ask many
  questions about why it is that conventional Western harmony has such a grip on our
  brains.
Jaki, Stanley L. Brain, Mind, and Computers. South Bend. Ind.: Gateway Editions, 1969.
  Paperback. A polemic book whose every page exudes contempt for the computational
  paradigm for understanding the mind. Nonetheless it is interesting to ponder the points
  he brings up.
• Jauch, J. M. Are Quanta Real? Bloomington, Ind.: Indiana University Press, 1973. A
  delightful little book of dialogues, using three characters borrowed from Galileo, put in
  a modern setting. Not only are questions of quantum mechanics discussed, but also
  issues of pattern recognition, simplicity, brain processes, and philosophy of science
  enter. Most enjoyable and provocative.

Bibliography                                                                            750

• Jeffrey, Richard. Formal Logic: Its Scope and Limits. New York: McGraw-Hill, 1967.
  An easy-to-read elementary textbook whose last chapter is on Godel's and Church's
  Theorems. This book has quite a different approach from many logic texts, which
  makes it stand out.
• Jensen, Hans. Sign, Symbol, and Script. New York: G. P. Putnam's, 1969. A-or perhaps
  the-top-notch book on symbolic writing systems the world over, both of now and long
  ago. There is much beauty and mystery in this book-for instance, the undeciphered
  script of Easter Island.
Kalmar, Laszlo. "An Argument Against the Plausibility of Church's Thesis". In A.
  Heyting, ed. Constructivity in Mathematics: Proceedings of the Colloquium held at
  Amsterdam, 1957, North-Holland, 1959. An interesting article by perhaps the
  bestknown disbeliever in the Church-Turing Thesis.
* Kim, Scott E. "The Impossible Skew Quadrilateral: A Four-Dimensional Optical
  Illusion". In David Brisson, ed. Proceedings of the 1978 A.A.A.S. Sym osium on
  Hypergraphics: Visualizing Complex Relationships in Art and Science. Boulder, Colo.:
  Westview Press, 1978. What seems at first an inconceivably hard idea-an optical
  illusion for four-dimensional "people"-is gradually made crystal clear, in an amazing
  virtuoso presentation utilizing a long series of excellently executed diagrams. The form
  of this article is just as intriguing and unusual as its content: it is tripartite on many
  levels simultaneously. This article and my book developed in parallel and each
  stimulated the other.
Kleene, Stephen C. Introduction to Mathematical Logic. New York: John Wiley, 1967. A
  thorough, thoughtful text by an important figure in the subject. Very worthwhile. Each
  time I reread a passage, I find something new in it which had escaped me before.
. Introduction to Metamathematics. Princeton: D. Van Nostrand (1952). Classic work on
  mathematical logic; his textbook (above) is essentially an abridged version. Rigorous
  and complete, but oldish.
Kneebone G. J. Mathematical Logic and the Foundations of Mathematics. New York:
  Van Nostrand Reinhold, 1963. A solid book with much philosophical discussion of
  such topics as intuitionism, and the "reality" of the natural numbers, etc.
Koestler, Arthur. The Act of Creation, New York: Dell, 1966. Paperback. A wide-ranging
  and generally stimulating theory about how ideas are "bisociated" to yield novelty. Best
  to open it at random and read, rather than begin at the beginning.
Koestler, Arthur and J. R. Smythies, eds. Beyond Reductionism. Boston: Beacon Press,
  1969. Paperback. Proceedings of Yr conference whose participants all were of the
  opinion that biological systems cannot be explained reductionistically, and that there is
  something "emergent" about life. I am intrigued by books which seem wrong to me, yet
  in a hard-to-pin-down way.
** Kutbose, Gyomay. Zen Koans. Chicago: Regnerv, 1973. Paperback. One of the best
  collections of koans available. Attractively presented- An essential book for any Zen
  library. Kuffler, Stephen W. and John G. Nicholls. From Neuron to Brain. Sunderland,
  deals Mss.: Smauer Associates, 1976. Paperback. A book which, despite its title, deals
  mostly with microscopic processes in the brain, and quite little with the way people's
  thoughts come out of the tangled mess. The work of Hubel and Wiesel on visual
  systems is covered particularly well.



Bibliography                                                                            751

Lacey Huh, and Geoffrey Joseph. "What the Godel Formula Says". Mind 77 (1968)- 7~.
  A useful discussion of the meaning of the Godel formula, based on a strict separation of
  three levels: uninterpreted formal system, interpreted formal system, a
  metamathematics. Worth studying. Latos Imre. Proofs and Refutations. New York:
  Cambridge University Press, 1976. Paperback. A most entertaining book in dialogue
  form, discussing how concepts are formed in mathematics. Valuable not only to
  mathematicians, but also to people interested in thought processes.
** Lehninger, Albert. Biochemistry. New York: Worth Publishers, 1976. A wonderfully
  readable text, considering its technical level. In this book one can find many ways in
  which proteins and genes are tangled together. Well organized, and exciting.
** Lucas, J. R. "Minds, Machines, and Godel". Philosophy 36 (1961): 112. This article is
  reprinted in Anderson's Minds and Machines, and in Sayre and Crosson's The Modeling
  of Mind. A highly controversial and provocative article, it claims to show that the
  human brain cannot, in principle, be modeled by a computer program. The argument is
  based entirely on Godel's Incompleteness Theorem, and is a fascinating one. The prose
  is (to my mind)
incredibly infuriating-vet for that very reason, it makes humorous reading. "Satan
  Stultified: A Rejoinder to Paul Benacerraf". Monist 52 (1968): 145.
Anti-Benacerraf argument, written in hilariously learned style: at one point Lucas refers
  to Benacerraf as "self-stultifyingly eristic" (whatever that means). The Lucas-
  Benacerraf battle, like the Lucas-Good battle, offers much food for thought.
  . "Human and Machine Logic: A Rejoinder". British Journal for the Philosophy of
  Science 19 (1967): 155. An attempted refutation of Good's attempted refutation of
  Lucas' original article.
** MacGillavry, Caroline H. Symmetry Aspects of the Periodic Drawings of M. C.
  Escher. Utrecht: A. Oosthoek's Uitgevermaatschappij, 1965. A collection of tilings of
  the plane by Escher, with scientific commentary by a crystallographer. The source for
  some of my illustrations-e.g., the Ant Fugue and the Crab Canon. Reissued in 1976 in
  New York by Harry N. Abrams under the title Fantasy and Symmetry.
MacKay, Donald M. Information, Mechanism and Meaning. Cambridge, Mass.: M.I.T.
  Press, 1970. Paperback. A book about different measures of information, applicable in
  different situations; theoretical issues related to human perception and understanding;
  and the way in which conscious activity can arise from a mechanistic underpinning.
• Mandelbrot, Benoit. Fractals: Form, Chance, and Dimension. San Francisco: W. H.
  Freeman, 1977. A rarity: a picture book of sophisticated contemporary research ideas in
  mathematics. Here, it concerns recursively defined curves and shapes, whose
  dimensionality is not a whole number. Amazingly, Mandelbrot shows their relevance to
  practically every branch of science.
• McCarthy, John. "Ascribing Mental Qualities to Machines". To appear in Martin
  Ringle, ed. Philosophical Perspectives in Artificial Intelligence. New York: Humanities
  Press, 1979. A penetrating article about the circumstances under which it would make
  sense to say that a machine had beliefs, desires, intentions, consciousness, or free will.
  It is interesting to compare this article with the book by Griffin.
Meschkowski, Herbert. Non-Euclidean Geometry. New York: Academic Press, 1964.
  Paperback. A short book with good historical commentary.



Bibliography                                                                            752

Meyer, can. ''Essai d'application de certains modeles cybernetiq ues it la coordina\_ Lion
  c1Jei les insectes sociaux". Insertes Sociaux XI11, no. 2 1966 : o which ch aws some
  parallels between the neural organization in the( brain, and the organiz ae lion of an ant
  colon\%.
Meyer, Leonard B. Emotion and Meaning in Music. Chicago: University of' Chicago
  Press, 1956. Paperback. A book which attempts to use ideas of Gestalt svcholo theory
  of perception to explain why musical structure is as it is. One of the more unusual the
  books on music and mind.
. Music, The Arts, and Ideas. Chicago: University of Chicago Press, 1967. Paperback. A
  thoughtful analysis of mental processes involved in listening to music, and of
  hierarchical structures in music. The author compares modern trends in music with Zen
  Buddhism.
Miller, G. A. and P. N. joint son-Laird. Language and Perception- Cambridge, Mass.:
  Harvard University Press, Belknap Press, 1976. A fascinating compendium of linguistic
  facts and theories, hearing on Whorl 's hypothesis that language is the same as
  worldviesv. A typical example is the discussion of the weird "mother-in-law" language
  of the Dyirbal people of Northern Queensland: a separate language used only for
  speaking to one's mother-in-law.
Minsky, Marvin L. "Matter, Mind, and Models". In Marvin L. Minskv, ed. Son antic
  Information Processing. Cambridge, Mass.: M.I.T. Press, 1968. Though merely a few
  pages long, this article implies a whole philosophy of consciousness and machine
  intelligence. It is a memorable piece of writing by one of the deepest thinkers in the
  field.
Minsky, Marvin L., and Seymour Papert\_ .92-t cial Intelligence Progress Report.
  Cambridge, Mass.: M.I.T. Artificial Intelligence Laboratory, Al Memo 252, 1972. A
  survey of all the work in Artificial Intelligence done at M.I.T. up to 1972, relating it to
  psychology and epistemology. Could serve excellently as an introduction to Al.
Monod, Jacques. Chance and Necessity. New York: Random House, Vintage Books,
  1971. Paperback. An extremel fertile mind writing in an idiosyncratic way about
  fascinating questions, such as how life is constructed out of non-life: how evolution,
  seeming\_ to violate the second law of thermodynamics, is actually dependent on it. The
  book excited me deeply.
* Morrison, Philip and Emil, eds. Charles Babbage and his Calculating Engines. Nesv
  York: Dover Publications, 1961. Paperback. A valuable source of information about the
  life of Babbage. A large fraction of Babbage's autobiography is reprinted here, along
  with several articles about Babbage's machines and his "Mechanical Notation".
Mvhill, John. "Some Philosophical Implications of Mathematical Logic". Review of
  Metaphysics 6 (1952): 165. An unusual discussion of ways in which Godel's Theorem
  and Church's Theorem are connected to psychology and epistemology. Ends up in a
  discussion of beauty and creativity.
Nagel, Ernest. The Structure of Science. New York: Harcourt, Brace, and World, 1961. A
  classic in the philosophy of science, featuring clear discussions of reductionism vs.
  holism, teleological vs. nonteleological explanations, etc.
Nagel, Ernest and James R. Newman. Godel's Proof. New York: New York University
  Press, 1958. Paperback. An enjoyable and exciting presentation, which was, in many
  ways, the inspiration for my own book.

Bibliography                                                                             753

* Nievergelt, Jurg, J. C. Farrar, and E. M. Reingold. Computer Approaches to Mathe,nat
  teal Problems. Englesyood Clif-f-s, N. J.: Prentice-Hall, 1974. An unusual collection of
  different types of problems which can be and have been attacked on computers-for
  instance, the "3n + I problem" (mentioned in my Rria with Diverse Variations) and
  other problems of number theory.
Pattee, Howard H., ed. Hierarchy Theory. New York: George Braziller, 1973. Paperback.
  Subtitled "The Challenge of Complex Systems". Contains a good article by Herbert
  Simon covering some of the same ideas as does my Chapter on "Levels of Description".
Peter, R6zsa. Recursive Functions. New York: Academic Press, 1967. A thorough
  discussion of primitive recursive functions, general recursive functions, partial
  recursive functions, the diagonal method, and many other fairly technical topics.
Quine, Willard Van Ornman. The Ways of Paradox, and Other Essays. New York:
  Random House, 1966. A collection of Quine's thoughts on many topics. The first essay
  deals with carious sorts of paradoxes, and their resolutions. In it, he introduces the
  operation I call "quining" in my book.
lishiRanganathan, S. R. Ramanujan, The Man and the Mathematician. London: Asia
  Pubng House, 1967. An occult-oriented biography of the Indian genius by an admirer.
  An odd but charming book.
Reichardt. Jasia. Cybernetics, Arts, and Ideas. Boston: New York Graphic Society,
197 1. A weird collection of ideas about computers and art, music, literature. Some of it
  is definitely off the deep end-hut some of it is not. Examples of the latter are the articles
  "A Chance for Art" by J. R. Pierce, and "Computerized Haiku" by Margaret
  Masterman.Renyi, Alfred. Dialogues on Mathematics. San Francisco: Holden-Day,
  1967. Paper hack. Three simple but stimulating dialogues involving classic characters
  in history, trying to get at the nature of mathematics. For the general public.
** Rips Paul. Zen Flesh, Zen Bones. New York: Doubleday, Anchor Books. Paperback.
  This book imparts very well the flavor of Zen-its antirational, antilanguage, antireduc
iionistic, basically holistic orientation.Rogers, Hartley. Theory of Recursive Functions
  and Effective Computability. New York:McGraw-Hill, 1967. A highly technical
  treatise, but a good one to learn from. Containsdiscussions of many intriguing problems
  in set theory and recursive function theory.
Rokeach, Milton. The Three Christs of' Ypsilanti. New York: Vintage Books,
  1964.paperback. A study of schizophrenia and the strange breeds of "consistency"
  which arise in the afflicted. A fascinating conflict between three men in a mental
  institution, all of whom imagined they were God, and how they dealt with being
  brought face to face for many months.
Rose, Steven. The Conscious Brain, updated ed. New York: Vintage Books, 1976.
Paperback. An excellent book-probably the best introduction to the study of the brain.
  Contains full discussions of the physical nature of the brain, as well as philosophical
  discussions on the nature\_ of mind, reductionism vs. holism, free will vs. determinism,
  etc. from a broad, intelligent, and humanistic viewpoint. Only his ideas on Al are way
  off.
Rosenblueth, Arturo. Mind and Brain: A Philosophy of Science. Cambridge,
  Mass.:M.I.T. Press, 1970. Paperback. A well written book by a brain researcher who
  deals with most of the deep problems concerning mind and brain.



Bibliography                                                                               754

* Sagan, Carl, ed. Communication with Extraterrestrial Intelligence. Cambridge, Mass.:
  M.I.T. Press, 1973. Paperback. Transcripts of a truly far-out conference, where a stellar
  group of scientists and others battle it out on this speculative issue.
Salmon, Wesley, ed. Zeno's Paradoxes. New York: Bobbs-Merrill, 1970. Paperback. A
  collection of articles on Zeno's ancient paradoxes, scrutinized under the light of modern
  set theory, quantum mechanics, and so on. Curious and thought-provoking,
  occasionally humorous.
Sanger. F., et al. "Nucleotide sequence of bacteriophage 16X174 DNA", Nature 265
  (Feb. 24, 1977). An exciting presentation of the first laying-bare ever of the full
  hereditary material of any organism. The surprise is the double-entendre: two proteins
  coded for in an overlapping way: almost too much to believe.
Sayre, Kenneth M., and Frederick J. Crosson. The Modeling of Mind: Computers and
  Intelligence. New York: Simon and Schuster, Clarion Books, 1963. A collection of
  philosophical comments on the idea of Artificial Intelligence by people from a wide
  range of disciplines. Contributors include Anatol Rapoport, Ludwig Wittgenstein,
  Donald Mackay, Michael Scriven, Gilbert Ryle, and others.
* Schank, Roger, and Kenneth Colby. Computer Models of Thought and Language. San
  Francisco: W. H. Freeman, 1973. A collection of articles on various approaches to the
  simulation of mental processes such as language-understanding, belief-systems,
  translation, and so forth. An important Al book, and many of the articles are not hard to
  read, even for the layman.
Schrodinger, Erwin. What is Life? \& Mind and Matter. New York: Cambridge University
  Press, 1967. Paperback. A famous book by a famous physicist (one of the main
  founders of quantum mechanics). Explores the physical basis of life and brain; then
  goes on to discuss consciousness in quite metaphysical terms. The first half, What is
  Life?, had considerable influence in the 1940's on the search for the carrier of genetic
  information.
Shepard, Roger N. "Circularity in judgments of Relative Pitch". Journal of the Acoustical
  Society of America 36, no. 12 (December 1964), pp. 2346-2353. The source of the
  amazing auditory illusion of "Shepard tones".
Simon, Herbert A. The Sciences of the Artificial. Cambridge, Mass.: M.I.T. Press, 1969.
  Paperback. An interesting book on understanding complex systems. The last chapter,
  entitled "The Architecture of Complexity", discusses problems of reductionism versus
  holism somewhat.
Smart, J. J. C. "Godel's Theorem, Church's Theorem, and Mechanism". Synthese 13
  (1961): 105. A well written article predating Lucas' 1961 article, but essentially arguing
  against it. One might conclude that you have to be Good and Smart, to argue against
  Lucas...
** Smullvan, Raymond. Theory of Formal Systems. Princeton, N. J.: Princeton
  University Press, 1961. Paperback. An advanced treatise, but one which begins with a
  beautiful discussion of formal systems, and proves a simple version of Godel's Theorem
  in an elegant way. Worthwhile for Chapter 1 alone. What Is the Name of This Book?
  Englewood Cliffs, N. J.: Prentice-Hall, 1978. A book of puzzles and fantasies on
  paradoxes, self-reference, and Godel's Theorem. Sounds like it will appeal to many of
  the same readers as my book. It appeared after mine was all written (with the exception
  of a certain entry in my bibliography).

Bibliography                                                                            755

Sommerhoff, Gerd. The Logic of the Living Brain. New York: John Wiley, 1974. A book
  which attempts to use knowledge of small-scale structures in the brain, in creating a
  theory of how the brain as a whole works.
Sperrv, Roger. "Mind, Brain, and Humanist Values". In John R. Platt, ed. .Vew Views on
  the Nature of Man. Chicago: University of Chicago Press, 1965, A pioneering
  neurophvsiologist here explains most vividly how he reconciles brain activity and
  consciousness.
* Steiner, George. After Babel: Aspects of Language and Translation. New York: Oxford
  University Press, 1975. Paperback. A book by a scholar in linguistics about the deep
  problems of translation and understanding of language by humans. Although Al is
  hardly discussed, the tone is that to program a computer to understand a novel or a
  poem is out of the question. A well written, thought-provoking-sometimes infuriating-
  book.
Stenesh, J. Dictionary of Biochemistry. New York: John Wiley, Wiley-Interscience,
  1975. For me, a useful companion to technical books on molecular biology.
** Stent, Gunther. "Explicit and Implicit Semantic Content of the Genetic Information".
  In The Centrality of Science and Absolute Values, Vol. 1. Proceedings of the 4th
  International Conference on the Unity of the Sciences, New York, 1975. Amazingly
  enough, this article is in the proceedings of a conference organized by the nowinfamous
  Rev. Sun Mvung Moon. Despite this, the article is excellent. It is about whether a
  genotype can be said, in any operational sense, to contain "all" the information about its
  phenotype. In other words, it is about the location of meaning in the genotype.
. Molecular Genetics: A Historical Narrative. San Francisco: W. H. Freeman, 1971. Stent
  has a broad, humanistic viewpoint, and conveys ideas in their historical perspective. An
  unusual text on molecular biology.
Suppes, Patrick. Introduction to Logic. New York: Van Nostrand Reinhold, 1957. A
  standard text, with clear presentations of both the Propositional Calculus and the
  Predicate Calculus. My Propositional Calculus stems mainly from here.
Sussman, Gerald Jay. A Computer Model of Skill Acquisition. New York: American
  Elsevier, 1975. Paperback. A theory of programs which understand the task of
  programming a computer. The questions of how to break the task into parts, and of how
  the different parts of such a program should interact, are discussed in detail.
** Tanenbaum, Andrew S. Structured Computer Organization. Englewood Cliffs, N. J.:
  Prentice-Hall, 1976. Excellent: a straightforward, extremely well written account of the
  many levels which are present in modern computer systems. It covers
  microprogramming languages, machine languages, assembly languages, operating
  systems, and many other topics. Has a good, partially annotated, bibliography.
Tarski, Alfred. Logic, Semantics, Metamathematics. Papers from 1923 to 1938.
  Translated by J. H. Lot, New York: Oxford University Press, 1956. Sets forth Tarski's
  ideas about truth, and the relationship between language and the world it represents.
  These ideas are still having repercussions in the problem of knowledge representation
  in Artificial Intelligence.
Taube, Mortimer. Computers and Common Sense. New York: McGraw-Hill, 1961.
  Paperback. Perhaps the first tirade against the modern concept of Artificial Intelligence.
  Annoying.



Bibliography                                                                            756

Tietze, Heinrich. Famous Problems of Mathematics. Baltimore: Graylock Press, 1965. A
  book on famous problems, written in a very personal and erudite style. Good
  illustrations and historical material.
Trakhtenbrot, V. Algorithms and Computing Machines. Heath. Paperback. A discussion
  of theoretical issues involving computers, particularly unsohable problems such as the
  halting problem, and the word-equisalence problem. Short. which is nice.
Turing, Sara. Alan M. Turing. Cambridge, U. K.: WW'. Hefter \& Sons, 1959. A
  biography of the great computer pion4r. A mother's work of love.
* Ulam, Stanislaw. Adventures of a Mathematician. New York: Charles Scribner's, 1976.
  An autobiography written by a sixty-five-year old man who writes as if he were still
  twenty and drunk in love with mathematics. Chock-full of gossip about who thought
  who was the best, and who envied whom, etc. Not only fun, but serious.
,A'aitson, J. D. The Molecular Biology of the Gene, 3rd edition. Menlo Park, Calif.: W.
  A. Benjamin, 1976. A good book but not nearly as well organized as Lehninger's, in my
  opinion. Still almost every page has something interesting on it.
Webb, Judson. "Metamathematics and the Philosophy of Mind". Philosophy of Science
  35 (1968): 156. A detailed and rigorous argument against Lucas, which contains this
  conclusion: "My overall position in the present paper may be stated by saying that the
  mind-machine-GOdel problem cannot be coherently treated until the constructivity
  problem in the foundations of mathematics is clarified."
Weiss, Paul. "One Plus One Does Not Equal Two". In G. C. Quarton, T. Melnechuk, and
  F. O. Schmitt, eds. The Neurosciences:.4 Study Program. New York: Rockefeller
  University Press, 1967. An article trying to reconcile holism and reductionism, but a
  good bit too holism-oriented for my taste.
* Weizenbaum, Joseph. Computer Power and Human Reason. San Francisco: W. H.
  Freeman, Freeman, 19i6. Paperback. A provocative book by an early At worker who
  has come to the conclusion that much work in computer science, particularly in Al, is
  dangerous. Although I can agree with him on some of his criticisms, I think he goes too
  far. His sanctimonious reference to Al people as "artificial intelligentsia" is funny the
  first time, but becomes tiring after the dozenth time. Anyone interested in computers
  should read it.
Wheeler, William Morton. "The Ant-Colony as an Organism".Journal of Morphology
22, 2 (1911): 307-325. One of the foremost authorities of his time on insects gives a
famous statement about why an ant colony deserves the label "organism" as much as its
  parts do.
Whitely, C. H. "Minds, Machines, and Godel: A Reply to Mr Lucas". Philosophy 37
  (1962): 61. A simple but potent reply to Lucas' argument. Wilder, Raymond. An
  Introduction to the Foundations of :Mathematics. New York: John Wiley, 1952. A good
  general overview, putting into perspective the important ideas of the past century.
* Wilson, Edward O. The Insect Societies. Cambridge, Mass.: Harvard University Press,
  elknap Press, 1971. Paperback. The authoritative book on collective behavior of insects.
  Although it is detailed, it is still readable, and discusses many fascinating ideas. It has
  excellent illustrations, and a giant (although regrettably not annotated) bibliography.
Winograd, Terry. Five Lectures on Artificial Intelligence. AI Memo 246. Stanford, Calif.:
  Stanford University Artificial Intelligence Laboratory, 1974. Paperback. A description
  of fundamental problems in At and new ideas for attacking them, by one of the

Bibliography                                                                             757

  important contemporary workers in the field.. Language as a Cognitive Process.
  Reading, Mass.: Addison-Wesley (forthcoming). From what I have seen of the
  manuscript, this will be a most exciting book, dealing with language in its full
  complexity as no other book ever has. . Understanding Natural Language. New York:
  Academic Press, 1972. A detailed discussion of one particular program which is
  remarkably "smart", in a limited world. The book shows how language cannot be
  separated from a general understanding of the world, and suggests directions to go in, in
  writing programs which can use language in the way that people do. An important
  contribution; many ideas can be stimulated by a reading of this book.
. "On some contested suppositions of generative linguistics about the scientific study of
  language", Cognition 4:6. A droll rebuttal to a head-on attack on Artificial Intelligence
  by some doctrinaire linguists.
* Winston, Patrick. Artificial Intelligence. Reading, Mass.: Addison-Wesley, 1977. A
  strong, general presentation of many facets of At by a dedicated and influential young
  proponent. The first half is independent of programs; the second half is LISP-dependent
  and includes a good brief exposition of the language LISP. The book contains many
  pointers to present-day At literature. ed. The Psychology of Computer Vision. New
  York: McGraw-Hill. 1975. Silly title, but fine book. It contains articles on how to
  program computers to do visual recognition of objects. scenes. and so forth. The articles
  deal with all levels of the problem, from the detection of line segments to the general
  organization of knowledge. In particular, there is an article by Winston himself on a
  program he wrote which develops abstract concepts from concrete examples, and an
  article by Minsky on the nascent notion of "frames".
* Wooldridge, Dean. Mechanical Man-The Physical Basis of Intelligent Life. New York:
  NJ McGraw-Hill, 1968. Paperback. A thorough-going discussion of the relationship of
  mental phenomena to brain phenomena, written in clear language. Explores difficult
  philosophical concepts in novel was, shedding light on them by means of concrete
  examples.




Bibliography                                                                           758

